% 10-appendix.tex 中文对照版（译自 sections/10-appendix.tex；不加 \label）
\appendix

\section{安全控制动作（UCA-01--08）}

表~\ref{tab:uca} 复现实现中定义的闸门八条不安全控制动作（源字符串为中文，此处按英文版转述）。该表、两个触发词表与判定逻辑硬编码于 \texttt{engine/control.py:175--188} 与 \texttt{engine/control.py:191--232}；检查按代码顺序求值，UCA-01--07 中任一命中都会把判定置为 \texttt{block}。UCA-08 只做降级：它被记录在 \texttt{downgrades} 字段中，从不改变判定。

\begin{table}[t]
\centering\small
\caption{不安全控制动作及其响应。“Marker”是固定的五字符锚点哨兵，以 U+3016 开头；“fact block”是渲染后的事实块 \texttt{f}。}
\begin{tabular}{@{}p{0.08\textwidth}p{0.42\textwidth}p{0.34\textwidth}p{0.09\textwidth}@{}}
\toprule
ID & 定义（由代码表转述） & 检测条件 & 响应 \\
\midrule
UCA-01 & 注入内容覆盖宪法约束 & $p$ 中任一覆盖短语（\texttt{:207--208}） & block \\
UCA-02 & 宪法校验失败但仍继续注入 & $\neg\mathrm{frozen\_ok}$（\texttt{:209--210}） & block \\
UCA-03 & 拒绝场景中注入了风格/情感 & $\mathrm{refuse} \land \mathrm{style\_injected}$（\texttt{:211--212}） & block \\
UCA-04 & 记忆内容被当作当前事实断言（无来源） & $f$ 中任一无来源短语（\texttt{:213--214}） & block \\
UCA-05 & 预算截断挤掉宪法 & $\neg\mathrm{constitution\_injected}$（\texttt{:215--216}） & block \\
UCA-06 & 未解决冲突未被降维 & unresolved $\land \neg$fallback（\texttt{:217--218}） & block \\
UCA-07 & 相变期人格锚点缺失 & prefeed $\land$ marker $\not\subseteq p$（\texttt{:219--220}） & block \\
UCA-08 & 饱和后积分仍继续 & $\exists k: s_k$（\texttt{:221--223}） & 仅降级 \\
\bottomrule
\end{tabular}
\end{table}

\section{参数与 provenance}

表~\ref{tab:params} 列出第~\ref{sec:method} 节各机制涉及的每个数值参数。权威来源是 \texttt{rules/decay-params.yaml}；\texttt{engine/config.py} 是零依赖的影子副本，由 \texttt{tools/check\_config\_sync.py} 逐键交叉核对（66 键，0 失败）。\textbf{所有条目均为 \texttt{provenance: guess}}；measured 参数数量为零（\texttt{acceptance/provenance-candidates.json} 对每个条目的 \texttt{upgrade\_eligible} 均为 \texttt{false}）。

\begin{table}[t]
\centering\footnotesize
\caption{数值参数、权威规则文件行号、取值、provenance 与代码影子。“guess”是设计默认值，不是标定值。}
\begin{tabular}{@{}p{0.25\textwidth}p{0.42\textwidth}p{0.18\textwidth}p{0.11\textwidth}@{}}
\toprule
参数键 & 取值 & Provenance & 影子 \\
\midrule
\texttt{actr.d}, \texttt{actr.ref}, \texttt{actr.lo}, \texttt{actr.hi} & 0.5, 0.001 d, $-3.0$, $1.0$（\texttt{:217--221}） & guess & \texttt{config.py:90--95} \\
\texttt{memory.actrWeight} & 0.6（\texttt{:18}） & guess\textsuperscript{a} & \texttt{config.py:21} \\
\texttt{memory.decay.tau.*} & 7, 30, 90 d（\texttt{:39--48}） & guess & \texttt{config.py:15--19} \\
\texttt{control.integralKi} / \texttt{Cap} / \texttt{Leak} / \texttt{Deadband} & 0.15 / 0.4 / 0.05 / 0.02（\texttt{:224--233}） & guess & \texttt{config.py:97--100} \\
\texttt{control.phaseShiftRatio} / \texttt{HardRatio} & 0.70 / 0.85（\texttt{:228--233}） & guess & \texttt{config.py:101--102} \\
\texttt{control.trendThreshold} / \texttt{ratioMax} / \texttt{historyCapacity} & 0.02 / 1.5 / 64（\texttt{:230--233}） & guess & \texttt{config.py:104--106} \\
\texttt{compose.weights.fact} / \texttt{organization} / \texttt{style} / \texttt{retrieval} & 0.70 / 0.50 / 0.40 / 0.45（\texttt{:239--247}） & guess & \texttt{config.py:50--53} \\
\texttt{compose.weights.boost*} & 0.50 / 0.35 / 0.40（\texttt{:243--247}） & guess & \texttt{config.py:54--56} \\
\texttt{compose.weights.styleBoostThreshold} & 0.6（\texttt{:246}） & guess & \texttt{config.py:58} \\
\texttt{memory.retrieval.topK} & 3（\texttt{:20}） & guess & \texttt{config.py:22} \\
\texttt{memory.recallBudget.tokenBudgetDefault} / \texttt{maxReturned} & 900 / 3（\texttt{:70--71}） & guess\textsuperscript{a} & \texttt{config.py:25} \\
\texttt{memory.potency.archiveThreshold} & 0.15（\texttt{:50}） & guess\textsuperscript{a} & \texttt{config.py:20} \\
\texttt{compose.injectionTokenRatioMax} & 0.15（\texttt{:236}） & none\textsuperscript{b} & \texttt{config.py:45}（\texttt{tokenRatioMax}） \\
\texttt{metacog.emotionCapIntensity} & 0.6 & 无规则键\textsuperscript{c} & \texttt{config.py:30} \\
\texttt{ASSUMED\_CONTEXT\_TOKENS} & 8000 & 无规则键\textsuperscript{c} & \texttt{compose.py:36} \\
\bottomrule
\end{tabular}
\end{table}

\noindent\textsuperscript{a}~该块没有内联 provenance 键；文件头声明所有值均为 guess（\texttt{:4}），其所属小节也重复了 \texttt{provenance: guess}。
\textsuperscript{b}~规则键存在，但一致性工具未比对它（与影子副本存在键名不匹配）。
\textsuperscript{c}~硬编码在代码中；作为单源例外记录在案。
这三处例外是仅有的“数值不来自被校验的规则文件键”的位置。

\section{复现命令与断言清单}

以下所有 Python 命令都在准备本文时于参考实现上执行过。在 Windows 上必须设置该环境变量，因为有一个断言名包含 U+2212，默认控制台代码页无法编码；不设置时，运行会在打印阶段中止，而不是在断言处失败。

\begin{verbatim}
$env:PYTHONIOENCODING = "utf-8"
python tests/test_control.py                  # 161 passed, 0 failed
python tests/test_pipeline_integration.py     #  65 passed, 0 failed
python tests/test_constitution.py             #  26 passed, 0 failed
python tools/check_config_sync.py             #  66 ok, 0 fail, 0 skipped
python tools/gen_impl_status.py --check       #  OK (status matrix current)
python tools/run_learning.py --json           #  recall_at_3 = 1.0
node plugin/dsh-lingshu/host-sim.mjs --assert #  documented 188 assertions
node plugin/dsh-lingshu/control-sim.mjs --assert # documented 148 assertions
node plugin/dsh-lingshu/stress.mjs --assert   #  documented 45 assertions
\end{verbatim}

表~\ref{tab:assertions} 把已执行的 Python 断言归因到第~\ref{sec:method} 节的机制。计数是执行次数，不是调用点：循环内的一个调用点会每轮执行一次（例如 top-$k$ 槽位检查）。插件断言计数为项目记录值，本文未重新执行。

\begin{table}[t]
\centering\footnotesize
\caption{断言清单。“Run”是观测到的断言执行次数；“site”是测试函数起始的测试文件行号。}
\begin{tabular}{@{}p{0.36\textwidth}p{0.34\textwidth}p{0.10\textwidth}p{0.12\textwidth}@{}}
\toprule
机制 & 测试函数（文件） & Run & Site \\
\midrule
ACT-R base level & \texttt{test\_actr\_activation}（\texttt{test\_control.py}） & 15 & \texttt{:58} \\
ACT-R dual channel & \texttt{test\_actr\_dual\_channel}（\texttt{test\_control.py}） & 7 & \texttt{:103} \\
Top-$k$ slot contract & \texttt{test\_topk\_slot\_contract}（\texttt{test\_control.py}） & 9 & \texttt{:143} \\
Integral feedback & \texttt{test\_integral\_baseline}（\texttt{test\_control.py}） & 18 & \texttt{:177} \\
Phase-shift monitor & \texttt{test\_phase\_shift\_monitor}（\texttt{test\_control.py}） & 22 & \texttt{:258} \\
Requisite-variety guard & \texttt{test\_requisite\_variety\_guard}（\texttt{test\_control.py}） & 11 & \texttt{:342} \\
Pre-flight gate & \texttt{test\_preflight\_audit}（\texttt{test\_control.py}） & 33 & \texttt{:387} \\
Dynamic weights & \texttt{test\_compose\_weights}（\texttt{test\_control.py}） & 29 & \texttt{:484} \\
Single-source parameters & \texttt{test\_params\_single\_source}（\texttt{test\_control.py}） & 8 & \texttt{:623} \\
Impasse & \texttt{test\_impasse}（\texttt{test\_control.py}） & 7 & \texttt{:668} \\
Pipeline integration & 九个函数 & 65 & \texttt{test\_pipeline\_integration.py:46--284} \\
Constitution audit & 五个函数 & 26 & \texttt{test\_constitution.py} \\
\bottomrule
\end{tabular}
\end{table}

正文引用的闸门级测量由与混沌执行器相同的无 token 调用产生（\texttt{tools/run\_chaos.py:443--446}）：样本文本作为 \texttt{prompt} 参数传入，并带 \texttt{frozen\_ok=True} 与 \texttt{constitution\_injected=True}，从而隔离出 UCA-01。这是对闸门扫描面的单元级覆盖检查，不是端到端拦截测量。
